\section*{Limitations}
\label{sec:limitations}

CoMem is evaluated primarily on one 8B backbone with a lexical selector and an
English benchmark suite. Results on other Qwen scales and sparse MoEs support
portability but are not matched replications, and several large-model cells use
smaller sample counts. The flagship is one batch-8 training run. Two additional
adapters match its data budget but use effective batch 3, so the reported
seed-plus-effective-batch check is not a clean estimate of training-seed
variance.

The paper does not provide a same-backbone, same-hardware implementation of the
closest PIC, chunk-KV repair, or learned modular-cache systems. Their published
operating points use different models, tasks, training budgets, storage tiers,
and latency boundaries, so Table~\ref{tab:priorart} is taxonomic and our
experiments establish an internal depth endpoint rather than competitive
superiority. The fully matched causal evidence is a
$j{=}0\!\rightarrow\!12$ two-point comparison; separately trained depths form a
deployment curve, not one controlled multi-depth frontier.

The system assumes repeated queries over a relatively stable corpus. Its
residual store grows linearly and is much larger than raw text or token IDs;
for one-off queries, frequently edited documents, or low reuse, raw-text
retrieval is usually preferable. Selector and external-store latency still grow
with corpus size even though model-side Read is bounded. Reported crossover and
speed measurements depend on hardware, kernels, batch size, storage tier, and
the precise timing boundary. At 128k, the displayed break-even result covers
only one generated token; longer-generation values were not retained and should
not be inferred from the 32k grid.

Stored states are tied to the exact backbone, tokenizer, split depth, adapter,
and lower-layer implementation; model upgrades generally require rewriting the
store. Overlap-Write also weakens strict chunk independence because a local edit
can invalidate the following chunk's cached state. We do not evaluate residual
quantization, eviction, or multi-tenant update contention.

Quality is task dependent. Bounded top-$k$ retrieval imposes an evidence ceiling,
lexical selection transfers poorly to low-overlap or non-whitespace languages,
and residual readout can fail even when the evidence is present. Overlap-Write
is tested on a focused multikey diagnostic rather than the natural benchmark
suite, and its wall-clock Write overhead is not integrated into a repaired
end-to-end frontier. It is not a universal quality fix. Results beyond the
backbone's native context window are positional-extrapolation stress tests.

The equal-latency mixture combines synthetic and natural diagnostics with equal
cell weights. Under iterative BM25, raw replay wins; with a frozen general
dense retriever, the aggregate is tied and task-level signs reverse. This
selector dependence precludes a selector-independent deployment conclusion.

LoCoMo's primary semantic metric uses an undated mutable \texttt{gpt-4o}
endpoint. We retain the evaluation date and saved item-level parsed decisions,
and a 200-item DeepSeek-V3 audit preserves the ordering, but neither makes the
full result exactly reproducible under a fixed judge snapshot. The PG-19 audit
also found substantial overlap with InfiniteBench long-book support, prompting
removal of that comparison. Equivalent overlap audits were not completed for
all natural benchmarks, including NarrativeQA; natural-task results are scope
checks, not contamination-free generalization estimates.

Our probes and controls characterize readout compatibility; they do not locate
``understanding,'' factual storage, or causal necessity at a particular layer.
Long-form generation, code, multilingual and multimodal memory, dynamic
updates/deletions, stronger learned retrievers, and production-scale
multi-tenant scheduling remain open. The serving experiments report medians on
single-query harnesses rather than concurrent p95 throughput or tail latency.
